\section{Variational load inversion with coupled viscous fluids}
\label{sec:viscous-load-inverse}

\subsection{The load is a multiplier of a physical evolution constraint}
Retain three incompressible Newtonian phases with positive viscosities,
piecewise constant densities, constant interfacial tensions, and pinned
contact lines at stationary no-slip walls. The interfaces may depend on both
horizontal coordinates. This subsection first excludes acoustic bulk forcing,
streaming and thermal evolution to define a mechanically realizable ideal-load
test. Those effects are restored as separate terms below, not claimed absent
in the eventual apparatus. Let $G_h\bm U$ be the graph kinematic map:
$(G_h\bm U)_j=U_z-\nabla h_j\cdot\bm U_{xy}$ on interface $j$.
Use $\langle f,v\rangle_D=\sum_j\int_D f_jv_j\dd x\dd y$.
The velocity is continuous across both interfaces, divergence-free in each
phase, and vanishes at the walls. Define the viscous bilinear form
\[
 a_h(\bm U,\bm W)=\sum_{\ell=0}^2\int_{\Omega_\ell(h)}
           2\mu_\ell\bm D(\bm U):\bm D(\bm W)\dd V.
\]
It has units of power when both arguments are velocities. The dissipation
potential is half this quadratic form. A Rayleighian construction of this
kind underlies viscous variational models \cite{qian2006}; here the contact
line is fixed and no moving-contact-line law is introduced.

For a supplied graph velocity $v_d$ compatible with all sealed volumes,
consider the saddle functional
\begin{equation}
 \mathscr R=\tfrac12 a_h(\bm U,\bm U)
       +\langle E'_h,G_h\bm U\rangle_D
       -\sum_\ell\int_{\Omega_\ell}p_\ell\nabla\cdot\bm U\dd V
       -\langle f,G_h\bm U-v_d\rangle_D.
 \label{eq:viscous-rayleighian}
\end{equation}
Here $E'_h$ is the full capillary--gravity variation. Variation in $p$ and
$f$ enforces incompressibility and the desired interface speed. Variation in
$\bm U$ gives the three-fluid Stokes weak equation. Thus $f$ is the total
required normal load, not a guessed pressure amplitude. Tangential traction
continuity is a natural condition of this scalar normal-load problem. If a
desired evolution requires incompatible tangential or bulk forces, normal
pressure alone does not realize it.

On the admissible volume-preserving graph-velocity space define
\begin{equation}
 \mathcal D_h(v)=\min_{\nabla\cdot\bm U=0,\ G_h\bm U=v}
                      \tfrac12 a_h(\bm U,\bm U)
             =\tfrac12\langle v,\mathcal Z_hv\rangle_D.
 \label{eq:viscous-resistance}
\end{equation}
Existence and the operator domain require suitable trace regularity and
compatible boundary data. Reciprocity and nonnegative dissipation imply a
symmetric positive resistance $\mathcal Z_h$ on that domain. Its off-diagonal
blocks couple the two interfaces through the middle fluid. Where its inverse
exists, $\mathcal M_h=\mathcal Z_h^{-1}$ is the Stokes mobility. The load and
forward equations are respectively
\begin{equation}
 \boxed{f_{\rm req}=E'_h+\mathcal Z_hv_d-\lambda,
 \qquad \dot h=\mathcal M_h(f-E'_h).}
 \label{eq:viscous-load-and-forward}
\end{equation}
Pressure constants are understood modulo the volume constraints. Neither a
diagonal relaxation constant nor an independently moving front and back is
consistent with this construction. At $v_d=0$ the viscous term vanishes,
recovering exactly the earlier static load. Adding viscosity does not change
a quiescent equilibrium pressure requirement; it changes the approach to it.

\subsection{What changes when inertia and additional forcing are retained}
With material acceleration retained, the inverse weak equation for a desired
admissible fluid motion is
\begin{align}
 \langle f-E'_h,G_h\bm W\rangle_D
 &=\sum_\ell\int_{\Omega_\ell}\rho_\ell
       (\partial_t\bm U+\bm U\cdot\nabla\bm U)\cdot\bm W\dd V
       +a_h(\bm U,\bm W) \nonumber\\
 &\quad-\sum_\ell\int_{\Omega_\ell}p_\ell\nabla\cdot\bm W\dd V
       -\mathcal W_{\rm additional}(\bm W).
 \label{eq:inertial-inverse-weak}
\end{align}
The last term is the power pairing of any additional mean bulk and tangential
loads. A scalar $f$ exists only if the right-hand side factors through $G_h$:
it must vanish for every admissible test velocity with $G_h\bm W=0$ after the
bulk equations and pressure are solved. This is a necessary compatibility
condition, not automatic controllability of a prescribed shape history.
Hydrodynamic memory means that shape and its instantaneous velocity alone
need not determine the required load. If the endpoint is a streaming state,
convective and viscous stresses persist there. With heating, density and
surface tension change and Marangoni stresses appear; the constant-property
static map is no longer the complete required acoustic load.

\subsection{A forward test that does not prescribe its answer}
For the ideal-load diagnostic choose the fixed full-domain static map
$f^*=E'_{h^*}-\lambda$ and start the forward problem at flat compatible
interfaces. The target appears only in this fixed load, not in a time-varying
shape prescription. In the Stokes limit its loaded potential obeys
\[
 \frac{\dd\mathcal U}{\dd t}
 =-\langle E'_h-f^*,\mathcal M_h(E'_h-f^*)\rangle_D\leq0.
\]
For a finite-dimensional volume-constrained discretization with bounded,
positive mobility and a strongly convex loaded potential on an invariant
admissible set, convergence to its unique minimizer follows. A continuum or
large-deformation assertion needs the corresponding regularity, gap and
coercivity bounds; a computed time series is not their proof.

The first numerical implementation uses axisymmetric Taylor--Hood velocity
and pressure elements on a moving, interface-fitted three-layer mesh. Velocity
is continuous, pressure is continuous within each phase but allowed to jump
at the interfaces. Only one global pressure gauge is removed. Both interface
volume fluxes follow from the phasewise divergence constraints. The full
axisymmetric strain includes hoop strain. The graph kinematics are projected
into pinned volume-preserving spline coordinates; the transpose of the same
map transfers generalized loads into the fluid. This reciprocal pairing
preserves the discrete work identity. The mesh represents curved interfaces
by line segments and must be refined separately from the spline space.

This first implementation is creeping flow: inertia, non-axisymmetric flow,
streaming and heat transport are not implemented by it. Viscosities are
explicit scenario inputs, not measured values inferred from desired settling.
Report velocity-based Reynolds estimates and a viscous diffusion-time check
before using it for a physical formation claim. In particular an abrupt load
switch in the Stokes model has no resolved inertial startup. Time refinement,
fluid-grid refinement, independent analytic flow limits and energy checks are
required. A pressure-history or array design cannot be declared sufficient
merely because this restricted diagnostic reaches its endpoint.
